\section{Introduction}
\label{sec:intro}

The relentless scaling of deep neural networks has democratized fine-tuning, allowing practitioners to adapt pre-trained foundational models to a diverse range of task-specific domains~\cite{Dosovitskiy2021, Radford2021, Devlin2019}. However, maintaining and deploying a separate, full-sized model for every unique task is computationally prohibitive at scale. This challenge has catalyzed the field of \emph{model merging} (or parameter fusion), which seeks to blend multiple specialized neural network experts into a single unified multitask architecture without requiring expensive joint re-training~\cite{Wortsman2022, Ilharco2022, Yadav2023}.

Early approaches to model merging focused on static methods, computing a single set of fixed blending coefficients across all layers and tasks~\cite{Matena2022, Jin2022dataless, OFS2025}. While static model merging achieves remarkable compression, it is fundamentally bounded: the blended model must use the same static weights regardless of whether the incoming query is from a simple task (such as digit recognition) or a complex task (such as out-of-distribution natural scenes). 

To overcome this limitation, recent literature has shifted towards \emph{dynamic model merging}~\cite{AdaMerging2023, vance2025quantum}. Dynamic merging architectures deploy a lightweight router that inspects input representations on-the-fly, dynamically predicting sample-specific blending coefficients to construct specialized model parameters on a query-by-query basis. While the promise of dynamic merging is substantial, it has recently fallen victim to a trend of escalating structural complexity. 

\begin{figure}[t]
\vskip 0.2in
\begin{center}
\centerline{\includegraphics[width=0.9\columnwidth]{comparison_plot.png}}
\caption{Homogeneous test stream accuracy comparison on the 4-task vision benchmark. Our proposed \textbf{Robust Linear Routing (RLR)}, despite having only 768 trainable routing parameters and utilizing standard classical regularization, stabilizes dynamic parameter fusion, preventing the out-of-distribution collapse on SVHN that plagues unregularized routing.}
\label{fig:comparison}
\end{center}
\vskip -0.2in
\end{figure}

The prime example of this trend is \emph{Quantum Wavefunction Superposition Merging (QWS-Merge)}~\cite{vance2025quantum}, which claims that classical linear routing is structurally incapable of robust dynamic parameter blending. Specifically, Vance~\cite{vance2025quantum} reported that a standard classical Linear Router collapses catastrophically on challenging, high-variance domains like the Street View House Numbers (SVHN) dataset, achieving a mere $15.30\%$ classification accuracy. To resolve this, QWS-Merge introduces a highly complex, multi-stage, quantum-inspired framework—formulated around task weight wavefunctions, phase-basis projectors, and dynamic wave interference patterns. While QWS-Merge manages to recover from the SVHN collapse, it introduces severe mathematical obfuscation, large backpropagation graphs, and significant runtime overhead.

In this work, we approach this problem through the guiding principle of \textbf{Occam's razor}: \emph{if a complex method can be matched or outperformed by a simpler one, the simpler one is strictly superior.} We critically analyze the reported failure of classical linear routing and find that the SVHN collapse is not a fundamental structural limitation of linear gating. Instead, it is a classic, unaddressed manifestation of \emph{overfitting} and \emph{high-variance routing logits} on out-of-distribution (OOD) representation spaces. During calibration on a tiny, few-shot dataset, the unregularized router easily learns extreme, overconfident weights that saturate the softmax output. When presented with high-variance SVHN inputs, these saturated weights make extreme, near-deterministic routing decisions that collapse the representational space of the merged model.

To prove this hypothesis, we propose \textbf{Robust Linear Routing (RLR)}, an extremely simple, minimalist dynamic merging framework. RLR retains the elegant, classical linear gating layer (requiring a mere 768 parameters) but stabilizes its optimization using two classical, timeless techniques:
\begin{enumerate}
    \item \textbf{$L_2$ Weight Decay:} Adds a standard Frobenius norm regularization penalty to the router's weights during calibration, bounding logit magnitudes and preventing extreme routing boundaries.
    \item \textbf{Softmax Temperature Scaling:} Introduces a constant temperature divisor ($T \ge 1$) to soften routing outputs, acting as an implicit entropy regularizer that preserves a stable mixture of task experts.
\end{enumerate}

Our empirical results, evaluated on a Vision Transformer (ViT-Tiny) benchmark across four distinct domains (MNIST, FashionMNIST, CIFAR-10, and SVHN), are striking:
\begin{itemize}
    \item \textbf{Statistical Robustness of Classical Gating:} We show that when trained using standard practices, the classical unregularized Linear Router is already highly robust, achieving an outstanding $91.53\% \pm 0.41\%$ Joint Mean accuracy across 5 seeds. This demonstrates that RLR's weight decay and temperature scaling are empirically redundant in standard homogeneous environments where unregularized routing is already highly effective.
    \item \textbf{Specialized Stabilization for Heterogeneity and OOD Shifts:} Rather than a mandatory fix for homogeneous collapse, RLR serves as a specialized stabilizer. It constrains logit variance and prevents overconfident gating under out-of-distribution shifts, slightly trading off a minor margin of homogeneous peak performance to secure superior resilience to mixed-task heterogeneous test streams, preventing the batch-level coefficient collapse that plagues standard dynamic routers.
    \item \textbf{Extreme Parameter and Calibration Efficiency:} RLR achieves these results with absolute simplicity, requiring only 768 parameters, optimizing in under a second on a single GPU on 64 calibration samples, introducing zero runtime overhead, and maintaining an elegant, 100-line implementation.
\end{itemize}

Importantly, while our empirical validation is conducted on compact Vision Transformers to directly deconstruct prior work, the simplicity and classical nature of RLR are structurally designed to scale seamlessly to modern high-dimensional architectures like Large Language Models (LLMs). Specifically, RLR's regularized gating opens up highly promising LLM-compatible scaling pathways, including (i) utilizing sequence-level pooled routing signals to minimize token-level routing noise, (ii) dynamically blending lightweight low-rank adapters (LoRAs) to eliminate parameter footprint, and (iii) exploiting linear mode connectivity in high-dimensional weight spaces to preserve representation smoothness (see Section~\ref{sec:conclusion} for detailed scaling formulations).

By showing that standard, classical regularization is all we need to stabilize and protect dynamic model merging, we hope to redirect the field's focus away from over-engineered complexity and back toward elegant, robust, and mathematically transparent solutions.
